\chapter{Alternative Data Strategies}\label{ch:s1:alternative-data-strategies}

When satellite images of retailers' car parks first reached investors, those who bought them could trade ahead of the retailers' quarterly reports,
especially the bad ones; Katona, Painter, Patatoukas and Zeng found more informed short selling and lower liquidity around the reports of covered retailers,
and concluded that unequal access to such data increased information asymmetry without immediately improving price discovery. An alternative dataset is
worth most to its first buyers and less to each one after. On the synthetic market a panel that reads a tenth of the companies' earnings surprises ten days
early earns a Sharpe ratio of 4.4 when no one else has it, 1.9 when half the market does, and loses when nine tenths do; a firm that gets the same data one
day before everyone else keeps 4.3 whatever the diffusion. This chapter takes a dataset to a position and measures how its value decays as it spreads. The
build is \texttt{firm.altstrat}.

\section{From dataset to position}

\begin{definition}[Nowcast, alternative-data strategy]\label{def:s1:alternative-data-strategies:nowcast}
A \emph{nowcast}\index{nowcast} is an estimate of a quantity that will be reported later (sales, earnings, economic activity) made from data observed
before the report. An \emph{alternative-data strategy}\index{alternative-data strategy} trades a security on nowcasts built from data outside the usual
market and accounting sources (transactions, images, web traffic, text), ahead of the report that will reveal the quantity.
\end{definition}

The path from a dataset to a position has four steps. Map the data to the companies it measures (Book~7, chapter~12 discussed coverage and the
point-in-time trap of backfilled histories). Turn the readings into a \omterm{def:s1:alternative-data-strategies:nowcast}{nowcast} of the reported quantity: here, the earnings surprise. Turn the \omterm{def:s1:alternative-data-strategies:nowcast}{nowcast} into
an expected return: the announcement's price reaction per unit of surprise. And hold the position over the window in which the market learns what the data
already say. The chapter's panel covers a tenth of the synthetic companies (about 420 announcements a year), reads each surprise with noise twice its spread (a
correlation of 0.46), and is delivered ten trading days before the announcement. The \omterm{def:s1:alternative-data-strategies:nowcast}{nowcast} is the reading times a slope of 0.21 fitted on the first two
years; the book holds the \omterm{def:s1:alternative-data-strategies:nowcast}{nowcast} from the delivery's close to the announcement's close, scaled to a gross exposure of one.

\section{What has been documented to work}

The published record is a set of datasets, each tied to a quantity that prices learn slowly. Satellite coverage of retailers let sophisticated investors
target bad reports (Katona and co-authors). Crowdsourced employer reviews forecast earnings surprises: Green, Huang, Wen and Zhou found that firms whose
ratings improved significantly outperformed firms whose ratings fell, with the effect in current employees' views of career opportunities and senior
management. Economic links are an older example: Cohen and Frazzini found that prices do not promptly incorporate news about a firm's principal customers,
and that a strategy trading suppliers on their customers' returns earned monthly alphas of over 150 basis points; the synthetic market's supplier links
(Book~7, chapter~10) plant that effect. In each case the edge is the gap between when the data know and when the price does.

\section{Decay as the data spread}

\begin{definition}[Data decay]\label{def:s1:alternative-data-strategies:decay}
\emph{Data decay}\index{data decay} is the loss of a dataset's trading value as more investors acquire it and trade on it earlier, moving prices before the
report and leaving less of the reaction for each user; unlike the post-publication decay of a published anomaly, it is driven by the number of buyers of the
data.
\end{definition}

Zhu found the other side of it: the introduction of consumer-transaction and satellite data raised stock price informativeness by lowering the cost of
acquiring information. More informative prices are good for markets and bad for the dataset's value. The chapter plants the mechanism directly
(\cref{lst:s1:alternative-data-strategies:diffuse}): a share $\phi$ of the market buys the same panel and moves each price on the delivery day by $\phi$ of
the announcement's expected reaction given the reading; that part is then missing from the announcement.

\begin{center}
\small
\begin{tabular}{@{}lccccc@{}}
\toprule
share of the market with the data & 0 & 0.25 & 0.5 & 0.75 & 0.9\\
\midrule
Sharpe ratio after costs & 4.41 & 3.20 & 1.86 & 0.45 & $-0.43$\\
return a year after costs (gross 1) & 61.1\% & 42.6\% & 24.2\% & 5.7\% & $-5.4\%$\\
\midrule
with the data one day before the others: Sharpe ratio & 4.28 & 4.40 & 4.43 & 4.37 & 4.29\\
\bottomrule
\end{tabular}
\end{center}

The value falls almost linearly with the share of the market that has the data, and turns negative before everyone has it, because the book pays ten
basis points per unit traded on a one-way turnover of 36 times its gross a year. Timing is the other half. The same dataset delivered to the firm one day before the others is
worth a Sharpe ratio of 4.3 to 4.4 at every level of diffusion (\cref{fig:s1:alternative-data-strategies:decay}): the firm trades before the others' move
and holds through it. In the dataset market this is why exclusivity and delivery speed are priced, and why vendors' early clients earn what their later
clients cannot.

\begin{omfigure}
\begin{tikzpicture}
\begin{axis}[width=10cm,height=5.4cm,xlabel={\small share of the market with the same data},ylabel={\small Sharpe ratio after costs},
  tick label style={font=\footnotesize},xmin=0,xmax=0.95,ymin=-1,ymax=5,grid=major,grid style={gray!25},legend pos=south west,
  legend style={font=\scriptsize},table/col sep=comma]
\addplot[omThm,very thick,mark=*] table[x=phi,y=sr_net]{figdata/strategies-1/16-alternative-data-strategies/decay.csv};
\addplot[omDef,very thick,mark=square*] table[x=phi,y=sr_head]{figdata/strategies-1/16-alternative-data-strategies/decay.csv};
\addplot[black!40,dashed] coordinates {(0,0) (0.95,0)};
\legend{data on the same day as the others,data one day earlier}
\end{axis}
\end{tikzpicture}

\omcaption{The \omterm{def:s1:alternative-data-strategies:nowcast}{nowcast} book on the synthetic market as the alternative panel spreads: Sharpe ratio after costs against the share of the market that buys the
same panel, for a firm that receives it on the same day as the others and for one that receives it a day earlier. Data:
\texttt{s1\_altstrat.run}.}\label{fig:s1:alternative-data-strategies:decay}
\end{omfigure}

\section{Operating an alternative-data strategy}

A live \omterm{def:s1:alternative-data-strategies:nowcast}{alternative-data strategy} is as much an operation as a model. The data arrive late, change format, lose coverage, and are revised; each delivery
has to be checked against its own history before it moves a position (Book~7, chapter~29's workflow). The vendor's history must be point in time, or the
backtest trades on backfilled readings fitted to the outcome (Book~7, chapter~12's trial). The \omterm{def:s1:alternative-data-strategies:nowcast}{nowcast} needs re-estimation as the data change. And the
strategy's own decay has to be monitored: the chapter's model suggests watching how much of the expected reaction happens before the report (the
pre-report move per unit of \omterm{def:s1:alternative-data-strategies:nowcast}{nowcast}), which rises as the data spread and is observable without knowing who else buys the data.

\section{Strategy files}

\begin{strategyfile}{Revenue nowcast trade}\label{strat:s1:alternative-data-strategies:nowcast}
\sfield{Who pays you, and why} Investors without the data, who learn the quarter's results only at the report.
\sfield{Instruments and venues} Stocks of companies the panel measures (retailers, consumer brands).
\sfield{Signal} A \omterm{def:s1:alternative-data-strategies:nowcast}{nowcast} of revenue or earnings against consensus, from transactions or images.
\sfield{Sizing and execution} Position from delivery to report, proportional to the \omterm{def:s1:alternative-data-strategies:nowcast}{nowcast}; hedge the market and sector.
\sfield{Costs} Data licences, often the largest cost; trading around reports.
\sfield{How it dies} Diffusion of the same data to more buyers; changes in the panel's coverage.
\sfield{Horizon, capacity, infrastructure} Weeks; capacity limited by the covered names; a data-engineering team.
\sfield{Backtest honestly} Point-in-time deliveries, not backfilled history; the report dates as scheduled.
\sfield{Sources} Katona, Painter, Patatoukas and Zeng (2024); Zhu (2019); this chapter's simulation.
\end{strategyfile}

\begin{strategyfile}{Web-traffic signal}\label{strat:s1:alternative-data-strategies:web}
\sfield{Who pays you, and why} As for \omterm{def:s1:alternative-data-strategies:nowcast}{nowcasts}: online activity leads reported sales.
\sfield{Instruments and venues} Companies whose sales happen online.
\sfield{Signal} Visits, app usage or searches against their seasonal pattern.
\sfield{Sizing and execution} As for \omterm{def:s1:alternative-data-strategies:nowcast}{nowcasts}.
\sfield{Costs} Data licences; noisy mapping from activity to revenue.
\sfield{How it dies} Changes in how traffic is measured; diffusion.
\sfield{Horizon, capacity, infrastructure} Weeks; entity mapping from domains and apps to companies.
\sfield{Backtest honestly} Traffic as measured at the time, with the vendor's methodology changes marked.
\sfield{Sources} No performance figure verified; Green et al.\ (2019) for the related case of employee reviews.
\end{strategyfile}

\begin{strategyfile}{Sentiment signal}\label{strat:s1:alternative-data-strategies:sentiment}
\sfield{Who pays you, and why} Investors slow to aggregate dispersed opinions about a company.
\sfield{Instruments and venues} Stocks with enough reviews, posts or articles.
\sfield{Signal} Changes in crowdsourced ratings or text tone (Book~7, chapter~12).
\sfield{Sizing and execution} Monthly sorts, or tilts in a broader book.
\sfield{Costs} Low turnover.
\sfield{How it dies} Manipulated reviews; diffusion.
\sfield{Horizon, capacity, infrastructure} Months; text processing.
\sfield{Backtest honestly} Reviews timestamped when posted; deleted and edited reviews handled.
\sfield{Sources} Green, Huang, Wen and Zhou (2019): improving employer ratings significantly outperform declining ones.
\end{strategyfile}

\begin{strategyfile}{Supply-chain linkage signal}\label{strat:s1:alternative-data-strategies:supply}
\sfield{Who pays you, and why} Attention-constrained investors who do not follow news through economic links.
\sfield{Instruments and venues} Suppliers and their principal customers.
\sfield{Signal} The customers' recent returns or news, mapped to their suppliers.
\sfield{Sizing and execution} Long suppliers of customers with good news, short the others; monthly.
\sfield{Costs} Moderate turnover.
\sfield{How it dies} Better-known links; link data now widely available.
\sfield{Horizon, capacity, infrastructure} A month; a supply-chain database, point in time.
\sfield{Backtest honestly} Links as disclosed at each date; customer disclosures' timing.
\sfield{Sources} Cohen and Frazzini (2008): monthly alphas of over 150 basis points; Book~7, chapter~10.
\end{strategyfile}

\section{Tutorial: by the time it is sold}\label{tut:s1:alternative-data-strategies:tut}

\begin{tutorial}
\textbf{Goal.} Turn a synthetic panel into \omterm{def:s1:alternative-data-strategies:nowcast}{nowcast} positions and measure how their value falls as the data spread, and what a head start is worth.
\textbf{End state:} the table and \cref{fig:s1:alternative-data-strategies:decay}.
\begin{tutsteps}
\item \textbf{Diffusion and the book}: others trade the data on its delivery; the book holds from delivery to report.
\omcode{code/firm/altstrat/firm_altstrat.py}{33}{54}{Diffusion of the data, and the pre-event book.}\label{lst:s1:alternative-data-strategies:diffuse}
\item \textbf{Run} \texttt{nowcast\_quality()}, \texttt{run(phi)} for five shares and \texttt{run(phi, 1)} for the head start, and \texttt{fig\_altstrat.py}.
\end{tutsteps}
\textbf{What to change next.} Let the share with the data grow over the years and measure the book's rolling Sharpe ratio; estimate the diffusion from the
pre-report move per unit of \omterm{def:s1:alternative-data-strategies:nowcast}{nowcast}; combine two panels with different delivery lags.
\end{tutorial}

\section{Build: alternative-data strategies}\label{bld:s1:alternative-data-strategies:altstrat}

\begin{build}
\bfield{Purpose} \omterm{def:s1:alternative-data-strategies:nowcast}{Nowcasts} from a noisy panel, the diffusion of the same data through the market, and the pre-report book.
\bfield{Interface} \texttt{measure(surprise, noise, rng)}, \texttt{nowcast\_slope(readings, surprises)}, \texttt{diffuse(R, events, jump, phi, lead, reading,
shrink)}, \texttt{pre\_event\_book(events, nowcast, lead, T, N)}.
\bfield{Rules} \omterm{def:s1:alternative-data-strategies:nowcast}{Nowcast} slopes fitted on earlier events only; diffusion keeps each announcement's total move unchanged; positions from the delivery's close.
\bfield{Acceptance tests} \texttt{code/firm/altstrat/tests/}: the slope of a noisy reading; diffusion moves the planted share and keeps the total; the book's
weights by hand.
\bfield{Stretch} Diffusion growing over time; vendors with revised histories; several panels with different lags.
\end{build}

\begin{omsources}
\item C.~Zhu, ``Big data as a governance mechanism'', \emph{Review of Financial Studies} 32(5), 2019.
\item Z.~Katona, M.~O.~Painter, P.~N.~Patatoukas and J.~Zeng, ``On the capital market consequences of big data: evidence from outer space'', \emph{Journal
of Financial and Quantitative Analysis} 60(2), 2025 (online 2024).
\item L.~Cohen and A.~Frazzini, ``Economic links and predictable returns'', \emph{Journal of Finance} 63(4), 2008.
\item T.~C.~Green, R.~Huang, Q.~Wen and D.~Zhou, ``Crowdsourced employer reviews and stock returns'', \emph{Journal of Financial Economics} 134(1), 2019.
\end{omsources}

\section{Exercises}

\begin{exercise}[$\star$]\label{exo:s1:alternative-data-strategies:1}
A panel reads a surprise with noise twice the surprise's standard deviation. What is the reading's correlation with the surprise, and the best \omterm{def:s1:alternative-data-strategies:nowcast}{nowcast} slope?
\end{exercise}

\begin{exercise}[$\star$]\label{exo:s1:alternative-data-strategies:2}
A stock's daily specific volatility is 1.6\% and the announcement jump is three of them per unit of surprise. With a \omterm{def:s1:alternative-data-strategies:nowcast}{nowcast} slope of 0.2 and a reading of 2,
what reaction does the \omterm{def:s1:alternative-data-strategies:nowcast}{nowcast} expect, and how much of it moves on delivery if half the market has the data?
\end{exercise}

\begin{exercise}[$\star$]\label{exo:s1:alternative-data-strategies:3}
Why is \omterm{def:s1:alternative-data-strategies:decay}{data decay} different from the post-publication decay of an anomaly?
\end{exercise}

\begin{exercise}[$\star\star$]\label{exo:s1:alternative-data-strategies:4}
Why does the book lose before the whole market has the data?
\end{exercise}

\begin{exercise}[$\star\star$]\label{exo:s1:alternative-data-strategies:5}
Why is a one-day head start worth almost the whole strategy?
\end{exercise}

\begin{exercise}[$\star\star$]\label{exo:s1:alternative-data-strategies:6}
How could you estimate, from prices alone, how widely your data are used?
\end{exercise}

\begin{exercise}[$\star\star\star$]\label{exo:s1:alternative-data-strategies:7}
\emph{Coding.} Run \texttt{run(0.5, 0, 3.0)}, with the panel's noise three times the surprise's spread instead of two. What happens to the Sharpe ratio, and why?
\end{exercise}

\begin{exercise}[$\star\star\star$]\label{exo:s1:alternative-data-strategies:8}
\emph{Find the flaw.} ``The vendor's five-year history shows our \omterm{def:s1:alternative-data-strategies:nowcast}{nowcast} would have earned a Sharpe ratio of 3; we will pay the licence.''
\end{exercise}

\section{Problem: By the Time It Is Sold}

\begin{problem}[{Weekend problem --- a dataset's life}]\label{pb:s1:alternative-data-strategies:1}
The chapter's synthetic panel and the public record.

\medskip\noindent\textbf{Part I --- From data to position.}
\begin{enumerate}
\item Define a \omterm{def:s1:alternative-data-strategies:nowcast}{nowcast} and an \omterm{def:s1:alternative-data-strategies:nowcast}{alternative-data strategy}.
\item List the four steps from a dataset to a position.
\item Describe the chapter's panel and its \omterm{def:s1:alternative-data-strategies:nowcast}{nowcast}.
\item What does the book hold, and when?
\end{enumerate}

\medskip\noindent\textbf{Part II --- The record.}
\begin{enumerate}[resume]
\item What did Katona and co-authors find about satellite data?
\item What did Green and co-authors find about employer reviews?
\item What did Cohen and Frazzini find about economic links?
\item What did Zhu find about price informativeness?
\end{enumerate}

\medskip\noindent\textbf{Part III --- Decay.}
\begin{enumerate}[resume]
\item Define \omterm{def:s1:alternative-data-strategies:decay}{data decay} and describe how the chapter plants it.
\item Give the Sharpe ratios as the share with the data grows.
\item Why does the value fall almost linearly?
\item What does a head start do?
\end{enumerate}

\medskip\noindent\textbf{Part IV --- The verdict.}
\begin{enumerate}[resume]
\item State the \emph{named result}: the \omterm{def:s1:alternative-data-strategies:nowcast}{nowcast} strategy's Sharpe ratio as the share of informed capital using the data grows.
\item How would you monitor a dataset's decay?
\item What operational risks does the strategy carry?
\item How would you evaluate a vendor's history?
\item What should a data licence be worth to the firm?
\item Which strategy file decays fastest?
\item How do these strategies relate to chapter~7?
\item In one sentence: what does an alternative dataset sell?
\end{enumerate}
\end{problem}

\section{Interview questions}

\begin{interviewq}[$\star$ \iqroles{researcher}]\label{iq:s1:alternative-data-strategies:1}
How would you evaluate a new alternative dataset?
\end{interviewq}

\begin{interviewq}[$\star\star$ \iqroles{researcher}]\label{iq:s1:alternative-data-strategies:2}
What is backfill bias in vendor data, and how do you detect it?
\end{interviewq}

\begin{interviewq}[$\star\star$ \iqroles{researcher, trader}]\label{iq:s1:alternative-data-strategies:3}
Your alternative-data signal has weakened over two years. What might be happening?
\end{interviewq}

\begin{interviewq}[$\star\star$ \iqroles{developer}]\label{iq:s1:alternative-data-strategies:4}
Design the pipeline that takes a daily vendor delivery to a trade the same evening.
\end{interviewq}

\begin{interviewq}[$\star\star$ \iqroles{risk}]\label{iq:s1:alternative-data-strategies:5}
What are the legal and compliance risks of alternative data?
\end{interviewq}

\begin{interviewq}[$\star\star\star$ \iqroles{researcher}]\label{iq:s1:alternative-data-strategies:6}
A share $\phi$ of the market trades a signal on its delivery and moves the price by $\phi$ of the expected reaction. Derive the expected P\&L of a trader who
receives it at delivery, and of one who receives it $k$ days earlier and holds through.
\end{interviewq}
